\chapter{Methodology}

\section{Analysis Approach}
The methodology should combine literature review, acoustic reasoning, and a validation-oriented design workflow. The first step is to identify the dominant cabin noise components and the likely transmission paths. The second step is to select a control topology that is compatible with the cabin layout. The third step is to estimate the benefit, complexity, and risk of the proposed design.

\section{Evaluation Metrics}
A professional assessment should use metrics that are meaningful for both engineering and passenger comfort.
\begin{itemize}
    \item Residual sound pressure level in the target zone.
    \item Frequency-selective attenuation in the low-frequency band.
    \item Computational cost of the control update.
    \item Sensitivity to modelling errors and sensor placement.
    \item Integration impact on mass, volume, and maintenance effort.
\end{itemize}

\section{Validation Plan}
A clear validation plan can be built in three levels. First, a simplified analytical model is used to explain the control principle. Second, a simulation model evaluates the system under representative disturbance scenarios. Third, a hardware-oriented discussion checks whether the selected architecture can be integrated into a real cabin environment.

\section{Expected Report Outputs}
The final competition report should include the assumption set, the control block diagram, the evaluation criteria, and the rationale for choosing the final architecture. If data becomes available later, the same structure can be extended with simulation plots, frequency response results, and a comparison against alternative designs.
